\documentclass[11pt]{article}
\usepackage{mathblatt}
% Präambel des Setzers (werkzeuge/setzer-praeambel.tex, wortgleich); dazu \gteil (Muster T2U)
\makeatletter
\setlength{\mbpfbe}{0mm}% keine Punkte (pruefung.md § 5: nur beim Prüfstein)
\newlength{\gEin}\setlength{\gEin}{0pt}
\let\gSchrift\relax
\renewcommand{\pfkreuzzeile}[1]{\par\noindent\hangindent3.2em\hangafter1\hspace*{1.6em}$\square$\ #1\par}
\newcommand{\gkreuz}[1]{$\square$\ #1\hspace{2em}}
\newcommand{\gkreuzzeile}[2]{\par\noindent\hangindent2.2em\hangafter1\mbox{#1\,$\square$}\ #2\par}% Buchstabe, Option: eine je Zeile (neben Skizze)
% links, was man ansieht, rechts, was man tut; Spaltengrenze überall an derselben Stelle (gemeinsam.md „Satz“)
\newsavebox{\gbox}\newlength{\gLinks}
% v0.10 (Lehrer 13.10.): keine starre Halbierung – die linke Spalte ist so breit wie die Figur (+ 6 mm), die
% Antwort steht direkt daneben; ist die Figur breiter als 0,62 Zeile, steht die Antwort darunter
\newcommand{\gzwei}[2]{\par\smallskip\noindent\sbox{\gbox}{#1}%
  \setlength{\gLinks}{\dimexpr\wd\gbox+6mm\relax}%
  \ifdim\gLinks>0.62\linewidth
    \usebox{\gbox}\par\smallskip\noindent\begin{minipage}{\linewidth}#2\end{minipage}%
  \else
    \begin{minipage}[t]{\gLinks}\vspace{0pt}\usebox{\gbox}\end{minipage}%
    \begin{minipage}[t]{\dimexpr\linewidth-\gLinks\relax}\vspace{0pt}\raggedright #2\end{minipage}%
  \fi\par}
% v0.9 (12.10., pruefung.md § 5 „Marke mit B, Bedingungen rechts“): drittes Argument von pfaufg = Bedingungen
% der Prüfung (v0.10: \pfohnetr = TR durchgestrichen, „FOR“), klein grau im rechten Rand auf der Grundlinie der Nummer, Zeilen mit \\
\newcommand{\gbedrand}[2]{\ifblank{#2}{}{\rlap{\hspace*{\dimexpr#1+2mm\relax}%
  \parbox[t]{15mm}{\raggedright\scriptsize\leavevmode\color{mbgrau}\strut#2}}}}
% Teil eines Bündels (a, b …): Marke links im Rand, Bedingungen rechts im Rand, auf der Zeile des Teils
\newcommand{\gteilrand}[2]{\llap{\mbpfrandmarke{#1}\hspace*{\mbpfnr}}\gbedrand{\linewidth}{#2}}
% Teil eines Bündels: Buchstabe in eigener Spalte, Inhalt eingerückt (\gEin = Einzug, damit die Spaltengrenze von
% \gzwei an derselben Stelle bleibt)
\newenvironment{gbteil}[3]{\par\smallskip\noindent\gteilrand{#2}{#3}\makebox[\mbpfteil][l]{\textbf{#1}}%
  \begin{minipage}[t]{\dimexpr\linewidth-\mbpfteil\relax}\setlength{\gEin}{\mbpfteil}\raggedright\strut\ignorespaces}%
 {\end{minipage}\par}
% Grafik höchstens so breit wie die Spalte (Bündel nebeneinander, Koordinatensystem ohne Karo)
\newsavebox{\gpbox}
\newcommand{\gpassend}[1]{\sbox{\gpbox}{#1}\ifdim\wd\gpbox>\linewidth\typeout{GPASSEND: verkleinert \the\wd\gpbox\space auf \the\linewidth}\resizebox{\linewidth}{!}{\usebox{\gpbox}}\else\usebox{\gpbox}\fi}
% Skizzenfeld (v0.10, Lehrer 13.10.): Karo ab dem linken Rand über die Zeile, „Skizze“ klein grau oben links im Feld
% (wie \pfskizzenfeld der Vorlage); der rechte Rand bleibt den Bedingungen
\newcommand{\gskizze}[1]{\par\smallskip\noindent\begin{tikzpicture}
  \clip (0,0) rectangle ({\linewidth-1mm},{#1*10mm});
  \draw[black!16,line width=0.3pt,step=5mm] (0,0) grid ({\linewidth},{#1*10mm});
  \node[anchor=north west,inner sep=2pt,font=\scriptsize,text=mbgrau,fill=white] at (0,{#1*10mm}) {Skizze};
  \end{tikzpicture}\par}
\RenewDocumentEnvironment{pfaufg}{m m m}{%
  \begin{lrbox}{\mbaufgabebox}\begin{minipage}[t]{\linewidth}%
  \gSchrift\noindent\hspace*{\gEin}\mbpfrandmarke{#1}\gbedrand{\linewidth-\gEin}{#3}%
  \begin{minipage}[t]{\mbpfnr}\strut\textbf{#2}\end{minipage}%
  \begin{minipage}[t]{\dimexpr\linewidth-\gEin-\mbpfnr\relax}\raggedright\strut\ignorespaces}%
 {\end{minipage}%
  \end{minipage}\end{lrbox}%
  \setlength{\mbaufgabehoehe}{\dimexpr\ht\mbaufgabebox+\dp\mbaufgabebox\relax}%
  \ifdefempty{\gkopfnach}{\par\addvspace{5pt}\Needspace*{\mbaufgabehoehe}}{%
    % Stufenkopf, der mit dieser (ersten) Aufgabe zusammenbleibt: Platz für beide verlangen, dann den Kopf setzen
    % (12.10., Quadratische: ein Kopf stand allein unten, die Aufgabe mit Koordinatensystem auf der nächsten Seite)
    \par\Needspace*{\dimexpr\mbaufgabehoehe+40pt\relax}\gkopfnach\gdef\gkopfnach{}\par\addvspace{5pt}}%
  \noindent\usebox{\mbaufgabebox}\par\addvspace{5pt}}
% Stufenkopf vormerken: \gkopfvor{<Kopf>} – die nächste pfaufg setzt ihn, sobald beide zusammen Platz haben
\def\gkopfnach{}
\newcommand{\gkopfvor}[1]{\gdef\gkopfnach{#1}}
% Teilaufgabe mit Skizze links in gzwei: Buchstabe in eigener Spalte, darunter Text und Skizze
\newcommand{\gteil}[2]{\makebox[\mbpfteil][l]{\textbf{#1}}\parbox[t]{\dimexpr0.5\textwidth-\gEin-\mbpfnr-5mm-\mbpfteil\relax}{\raggedright #2}}
% Stufenkopf der Vorlage (mathblatt.sty 2026-10-12: Linie, Abstand, 13 pt, \pfjahre) ohne ihren Vorhalt von
% 14 Zeilen: der Setzer fordert vor dem Kopf 0,16 Seitenhöhe (Kopf + erste Aufgabe), mehr nicht
\setlength{\mbpfstufevorhalt}{0pt}
\makeatother
\newcommand{\karo}[1]{\par\smallskip\noindent\begin{tikzpicture}
  \clip (0,0) rectangle ({\linewidth-1mm},{#1*10mm});
  \draw[black!16,line width=0.3pt,step=5mm] (0,0) grid ({\linewidth},{#1*10mm});
  \end{tikzpicture}\par}
\newcommand{\kk}[1]{$\square$\,#1\hspace{0.9em}}
\newcommand{\linie}[1]{\,{\color{black!45}\rule[-1pt]{#1}{0.4pt}}\,}
\newcommand{\pkt}{\textperiodcentered{}}
% v0.14 (14.10.): Stufenkopf mit fünf Punkten (Option = Jahre für die Bestellkarte) oder „zuletzt 19“
\NewDocumentCommand{\gpunkte}{o m}{#2}
\newcommand{\gzuletzt}[1]{zuletzt #1}
% v0.7 (12.10.): Schülerfrage unter dem Stufenkopf; Ankreuzoption als Bild (Buchstabe und Kästchen darunter)
\newcommand{\gfrage}[1]{\par\nobreak\noindent\hspace*{\mbpfnr}{\small\itshape #1}\par\nobreak\addvspace{2pt}}
\newlength{\gbildbreite}\setlength{\gbildbreite}{0.32\linewidth}% Setzer setzt sie je Reihe (vier Bilder: 0,24)
\newcommand{\gbildoption}[3]{\begin{minipage}[t][#3][b]{\gbildbreite}\centering #2\par\smallskip\mbox{#1\,$\square$}\end{minipage}}
% Lineare (12.10.): Bild mit Buchstabe ohne Kästchen (Diagramm zuordnen); Aussage mit „wahr / falsch“ zum Ankreuzen
\newcommand{\gbildname}[3]{\begin{minipage}[t][#3][b]{\gbildbreite}\centering #2\par\smallskip\mbox{\textbf{#1}}\end{minipage}}
\newcommand{\gwahrfalsch}[2]{\par\noindent\hangindent2.2em\hangafter1\makebox[2.2em][l]{(#1)}#2\hfill\mbox{wahr $\square$\quad falsch $\square$}\par}
\begin{document}

\pfheftstil
\pflaufkopf[Abitur GK \pkt{} Prüfung]{Rechter Winkel \pkt{} Lösungen}{AG30}\pfseitevon
\pfheftkopf{Rechter Winkel}{Lösungen}
\pfstufenzeile{rechten Winkel nachweisen \pkt{} Normalenvektor nachweisen \pkt{} Parameter oder Punkt gesucht}

\Needspace*{0.16\textheight}\pfstufekopf{rechten Winkel nachweisen}{\gpunkte[26 25 24]{\textcolor{black!80}{$\bullet$}\,\textcolor{black!80}{$\bullet$}\,\textcolor{black!80}{$\bullet$}\,\textcolor{black!25}{$\circ$}\,\textcolor{black!25}{$\circ$}}}\special{pdf:dest (stufe-rechten-winkel-nachweisen) [@thispage /XYZ @xpos @ypos null]}
\begin{pfloesungb}{}
\lzb{1.}{}{rechter Winkel in A}{$\overrightarrow{AB}$ · $\overrightarrow{AC}$ = 4t $-$ 4t + 0 = 0 für jedes t $\Rightarrow$ rechter Winkel in A}
\end{pfloesungb}
\begin{pfloesungb}{}
\lzb{2.}{}{ja}{ja: Richtungsvektor von h $\left(\begin{smallmatrix}0\\1\\0\end{smallmatrix}\right)$; $\left(\begin{smallmatrix}-3\\0\\2\end{smallmatrix}\right)$ · $\left(\begin{smallmatrix}0\\1\\0\end{smallmatrix}\right)$ = 0 $\Rightarrow$ senkrecht}
\end{pfloesungb}
\begin{pfloesungb}{}
\lzb{3.}{}{rechter Winkel bei C}{$\overrightarrow{CD}$ = $\left(\begin{smallmatrix}-3\\3\\0\end{smallmatrix}\right)$, $\overrightarrow{CS}$ = $\left(\begin{smallmatrix}-3\\-3\\5\end{smallmatrix}\right)$ $\Rightarrow$ $\overrightarrow{CD}$ · $\overrightarrow{CS}$ = 9 $-$ 9 + 0 = 0 $\Rightarrow$ rechter Winkel bei C}
\end{pfloesungb}

\Needspace*{0.16\textheight}\pfstufekopf{Normalenvektor nachweisen}{\gpunkte[25 22]{\textcolor{black!25}{$\circ$}\,\textcolor{black!80}{$\bullet$}\,\textcolor{black!25}{$\circ$}\,\textcolor{black!25}{$\circ$}\,\textcolor{black!80}{$\bullet$}}}\special{pdf:dest (stufe-normalenvektor-nachweisen) [@thispage /XYZ @xpos @ypos null]}
\begin{pfloesungb}{}
\lzb{4.}{}{senkrecht auf E}{$\left(\begin{smallmatrix}-4\\3\\1\end{smallmatrix}\right)$ · $\left(\begin{smallmatrix}3\\4\\0\end{smallmatrix}\right)$ = $-$12 + 12 = 0 und $\left(\begin{smallmatrix}4\\-3\\0\end{smallmatrix}\right)$ · $\left(\begin{smallmatrix}3\\4\\0\end{smallmatrix}\right)$ = 12 $-$ 12 = 0 $\Rightarrow$ senkrecht zu beiden Spannvektoren $\Rightarrow$ senkrecht auf E}
\end{pfloesungb}
\begin{pfloesungb}{}
\lzb{5.}{}{Normalenvektor von L und N; 5x + 5y + 3z = 30}{(5|5|3) · CG = $-$15 + 0 + 15 = 0, (5|5|3) · CF = 0 $-$ 15 + 15 = 0 $\Rightarrow$ Normalenvektor von L und N; B einsetzen $\Rightarrow$ N: 5x + 5y + 3z = 30}
\end{pfloesungb}

\Needspace*{0.16\textheight}\pfstufekopf{Parameter oder Punkt gesucht}{\gpunkte[26 25 24 22]{\textcolor{black!80}{$\bullet$}\,\textcolor{black!80}{$\bullet$}\,\textcolor{black!80}{$\bullet$}\,\textcolor{black!25}{$\circ$}\,\textcolor{black!80}{$\bullet$}}}\special{pdf:dest (stufe-parameter-oder-punkt-gesucht) [@thispage /XYZ @xpos @ypos null]}
\begin{pfloesungb}{}
\lzb{6.}{}{s = 1,2}{$\left(\begin{smallmatrix}2\\-5\\0\end{smallmatrix}\right)$ · $\left(\begin{smallmatrix}3\\s\\1\end{smallmatrix}\right)$ = 6 $-$ 5s = 0 $\Rightarrow$ s = 1,2}
\end{pfloesungb}
\begin{pfloesungb}{}
\lzb{7.}{}{t = 7}{$\overrightarrow{QO}$ · $\overrightarrow{QP}$ = 20 + t² $-$ 69 = t² $-$ 49 = 0 $\Rightarrow$ t = $-$7 oder t = 7}
\end{pfloesungb}
\begin{pfloesungb}{}
\lzb{8.}{}{rechter Winkel bei B'}{B' hat die Koordinaten (3|y|0), weil sich bei der Verschiebung parallel zur y-Achse nur y ändert; (3|y|0) ist $\overrightarrow{AB'}$, (3|y $-$ 8|0) ist $\overrightarrow{CB'}$; Skalarprodukt null $\Rightarrow$ rechter Winkel bei B'}
\end{pfloesungb}
\begin{pfloesungb}{}
\lzb{9.}{}{C($\tfrac{10}{3}$|$-\tfrac{5}{3}$|5)}{AB = (6|3|$-$3), AC = (x $-$ 1|y + 2|5); I: 6x + 3y $-$ 15 = 0, II: x + 2y = 0 $\Rightarrow$ y = $-$5/3, x = $\tfrac{10}{3}$ $\Rightarrow$ C($\tfrac{10}{3}$|$-\tfrac{5}{3}$|5)}
\end{pfloesungb}
\end{document}